\section{AHSME 1975}
        \begin{enumerate}
        	\item (\textit{AHSME 1975}) The value of $\frac {1}{2 - \frac {1}{2 - \frac {1}{2 - \frac12}}}$ is 


 $\textbf{(A)}\ 3/4 \qquad  \textbf{(B)}\ 4/5 \qquad  \textbf{(C)}\ 5/6 \qquad  \textbf{(D)}\ 6/7 \qquad  \textbf{(E)}\ 6/5$


 
 \textbf{Answer: }B



\textbf{Solution}

Solution by e\_power\_pi\_times\_i

 
Calculating, we find that $\frac {1}{2 - \frac {1}{2 - \frac {1}{2 - \frac12}}} = \frac {1}{2 - \frac {1}{2 - \frac {2}{3}}} = \frac {1}{2 - \frac {3}{4}} = \frac {1}{\frac {5}{4}} = \boxed{\textbf{(B) } \dfrac{4}{5}}$.

 


	\item (\textit{AHSME 1975}) For which real values of m are the simultaneous equations 

 \begin{align*}y &= mx + 3 \\  y& = (2m - 1)x + 4\end{align*}
 
satisfied by at least one pair of real numbers $(x,y)$? 


 $\textbf{(A)}\ \text{all }m\qquad \textbf{(B)}\ \text{all }m\neq 0\qquad \textbf{(C)}\ \text{all }m\neq 1/2\qquad \textbf{(D)}\ \text{all }m\neq 1\qquad \\ \textbf{(E)}\ \text{no values of }m$


 
 \textbf{Answer: }D



\textbf{Solution 1}

Setting the two expressions for $y$ equal to each other yields $mx+3 = (2m-1)x+4$, which simplifies to $(m-1)x+1 = 0$. Therefore $x = \dfrac{1}{1-m}$, so $x$ is a real number if and only if $\boxed{\textbf{(D) }m\neq 1}$.

 -e\_power\_pi\_times\_i

 



\textbf{Solution 2}

Any two lines with different slopes intersect.  So if $m \neq 2m - 1$, that is, $m \neq 1$, the lines intersect and thus the system has one solution.

 If $m = 1$, however, the lines will have equations $y = x + 3$ and $y = x + 4$.  These lines have the same slope but different $y$-intercepts, so they are parallel and thus the system has no solutions.  

 Therefore, the system has at least one solution if and only if $\boxed{\textbf{(D) }m\neq 1}$.

 -j314andrews

 


	\item (\textit{AHSME 1975}) Which of the following inequalities are satisfied for all real numbers $a, b, c, x, y, z$ which satisfy the conditions $x < a, y < b$, and $z < c$? 


 $\text{I}. \ xy + yz + zx < ab + bc + ca \\ \text{II}. \ x^2 + y^2 + z^2 < a^2 + b^2 + c^2 \\ \text{III}. \ xyz < abc$


 $\textbf{(A)}\ \text{None are satisfied.} \qquad  \textbf{(B)}\ \text{I only} \qquad  \textbf{(C)}\ \text{II only} \qquad  \textbf{(D)}\ \text{III only} \qquad  \textbf{(E)}\ \text{All are satisfied.}$


 
 \textbf{Answer: }A



\textbf{Solution}

Solution by e\_power\_pi\_times\_i

 
Notice if $a$, $b$, and $c$ are $0$, then we can find $x$, $y$, and $z$ to disprove $\text{I}$ and $\text{II}$. For example, if $(a, b, c, x, y, z) = (0, 0, 0, -1, -1, -1)$, then $\text{I}$ and $\text{II}$ are disproved. If $(a, b, c, x, y, z) = (0, 1, 2, -1, -1, 1)$, then $\text{III}$ is disproved. Therefore the answer is $\boxed{\textbf{(A) } \text{None are satisfied}}$.

 


	\item (\textit{AHSME 1975}) If the side of one square is the diagonal of a second square, what is the ratio of the area of the first square to the area of the second? 


 $\textbf{(A)}\ 2 \qquad  \textbf{(B)}\ \sqrt2 \qquad  \textbf{(C)}\ 1/2 \qquad  \textbf{(D)}\ 2\sqrt2 \qquad  \textbf{(E)}\ 4$


 
 \textbf{Answer: }A



\textbf{Solution}

Solution by e\_power\_pi\_times\_i

 
Denote the side of one square as $s$. Then the diagonal of the second square is $s$, so the side of the second square is $\dfrac{s\sqrt{2}}{2}$. The area of the second square is $\dfrac{1}{2}s^2$, so the ratio of the areas is $\dfrac{s^2}{\dfrac{1}{2}s^2} = \boxed{\textbf{(A) } 2}$.

 


	\item (\textit{AHSME 1975}) The polynomial $(x+y)^9$ is expanded in decreasing powers of $x$. The second and third terms have equal values 
when evaluated at $x=p$ and $y=q$, where $p$ and $q$ are positive numbers whose sum is one. What is the value of $p$? 


 $\textbf{(A)}\ 1/5 \qquad  \textbf{(B)}\ 4/5 \qquad  \textbf{(C)}\ 1/4 \qquad  \textbf{(D)}\ 3/4 \qquad  \textbf{(E)}\ 8/9$


 
 \textbf{Answer: }B



\textbf{Solution}

Solution by e\_power\_pi\_times\_i

 
The second and third term of $(x+y)^9$ is $9x^8y$ and $36x^7y^2$, respectively. For them to be equal when $x = p$, $\dfrac{p}{4} = y$. For them to be equal when $y = q$, $x = 4q$. Then $p+\dfrac{p}{4} = q+4q$, so $\dfrac{5p}{4} = 5q$, which simplifies to $p = 4q$. Since $p+q = 1$, $p = \boxed{\textbf{(B) } 4/5}$.

 


	\item (\textit{AHSME 1975}) The sum of the first eighty positive odd integers subtracted from the sum of the first eighty positive even integers is 


 $\textbf{(A)}\ 0 \qquad  \textbf{(B)}\ 20 \qquad  \textbf{(C)}\ 40 \qquad  \textbf{(D)}\ 60 \qquad  \textbf{(E)}\ 80$


 
 \textbf{Answer: }E



\textbf{Solution}

Solution by e\_power\_pi\_times\_i

 
When the $n$th odd positive integer is subtracted from the $n$th even positive integer, the result is $1$. Therefore the sum of the first eighty positive odd integers subtracted from the sum of the first eighty positive even integers is $80\cdot1 = \boxed{\textbf{(E) } 80}$.

 


	\item (\textit{AHSME 1975}) For which non-zero real numbers $x$ is $\frac{|x-|x|\-|}{x}$ a positive integer? 


 $\textbf{(A)}\ \text{for negative } x \text{ only} \qquad \\ \textbf{(B)}\ \text{for positive } x \text{ only} \qquad \\ \textbf{(C)}\ \text{only for } x \text{ an even integer} \qquad \\ \textbf{(D)}\ \text{for all non-zero real numbers } x \\ \textbf{(E)}\ \text{for no non-zero real numbers } x$


 
 \textbf{Answer: }E



\textbf{Solution}

Solution by e\_power\_pi\_times\_i

 
Notice that if $x$ is negative, then the whole thing would amount to a negative number. Also notice that if $x$ is positive, then $|x-|x|\-|$ would be $0$, hence the whole thing would amount to $0$. Therefore, $\frac{|x-|x|\-|}{x}$ is positive $\boxed{\textbf{(E)}\ \text{for no non-zero real numbers } x}$.

 


	\item (\textit{AHSME 1975}) If the statement "All shirts in this store are on sale." is false, then which of the following statements must be true? 


 I. All shirts in this store are at non-sale prices. 


 II. There is some shirt in this store not on sale. 


 III. No shirt in this store is on sale. 


 IV. Not all shirts in this store are on sale. 


 

$\textbf{(A)}\ \text{II only} \qquad  \textbf{(B)}\ \text{IV only} \qquad  \textbf{(C)}\ \text{I and III only} \qquad  \textbf{(D)}\ \text{II and IV only} \qquad  \textbf{(E)}\ \text{I, II and IV only}$


 
 \textbf{Answer: }D



\textbf{Solution}

Solution by e\_power\_pi\_times\_i

 
If not all the shirts in the store are on sale, then the answer is $\boxed{\textbf{(D) } \text{II and IV only}}$.

 


	\item (\textit{AHSME 1975}) Let $a_1, a_2, \ldots$ and $b_1, b_2, \ldots$ be arithmetic progressions such that $a_1 = 25, b_1 = 75$, and $a_{100} + b_{100} = 100$. 
Find the sum of the first hundred terms of the progression $a_1 + b_1, a_2 + b_2, \ldots$


 $\textbf{(A)}\ 0 \qquad  \textbf{(B)}\ 100 \qquad  \textbf{(C)}\ 10,000 \qquad  \textbf{(D)}\ 505,000 \qquad \\ \textbf{(E)}\ \text{not enough information given to solve the problem}$


 
 \textbf{Answer: }C



\textbf{Solution}

Notice that $a_{100}$ and $b_{100}$ are $25+99k_1$ and $75+99k_2$, respectively. Therefore $k_2 = -k_1$. Now notice that $a_n + b_n = 25+k_1(n-1)+75+k_2(n-1) = 100+k_1(n-1)-k_1(n-1) = 100$. The sum of the first $100$ terms is $100\cdot100 = \boxed{\textbf{(C) } 10,000}$.

 


	\item (\textit{AHSME 1975}) The sum of the digits in base ten of $(10^{4n^2+8}+1)^2$, where $n$ is a positive integer, is 


 $\textbf{(A)}\ 4 \qquad  \textbf{(B)}\ 4n \qquad  \textbf{(C)}\ 2+2n \qquad  \textbf{(D)}\ 4n^2 \qquad  \textbf{(E)}\ n^2+n+2$


 
 \textbf{Answer: }A



\textbf{Solution}

We see that the result of this expression will always be in the form $(100\text{ some number of zeros }001)^2.$ Multiplying these together yields: 
 \[110\text{ some number of zeros }011.\] This works because of the way they are multiplied. Therefore, the answer is $\boxed{(A) 4}$.

 


	\item (\textit{AHSME 1975}) Let $P$ be an interior point of circle $K$ other than the center of $K$. Form all chords of $K$ which pass through $P$, 
and determine their midpoints. The locus of these midpoints is 


 $\textbf{(A)}\ \text{a circle with one point deleted} \qquad \\ \textbf{(B)}\ \text{a circle if the distance from } P \text{ to the center of } K \text{ is less than} \\ \text{one half the radius of } K \text{; otherwise a circular arc of less than} 360^{\circ}\qquad \\ \textbf{(C)}\ \text{a semicircle with one point deleted} \qquad \\ \textbf{(D)}\ \text{a semicircle} \qquad  \textbf{(E)}\ \text{a circle}$


 
 \textbf{Answer: }E



\textbf{Solution 1}

It is (E) a circle.

 The midpoints of the chord seems to be on the circle. For example, draw a chord passing through P. It is on the circle PK, where PK is the diameter.To prove this, we can use the fact thatThe perpendicular bisector of any chord of a circle passes through the center of the circle.

 "Hence, the perpendicular through A passes through K. Since angle PAK=90, angle PAK is inscribed in the semicircular arc PK"

 (Volume 2 , Chapter 20)

 


	\item (\textit{AHSME 1975}) If $a \neq b, a^3 - b^3 = 19x^3$, and $a-b = x$, which of the following conclusions is correct? 


 $\textbf{(A)}\ a=3x \qquad  \textbf{(B)}\ a=3x \text{ or } a = -2x \qquad  \textbf{(C)}\ a=-3x \text{ or } a = 2x \qquad \\ \textbf{(D)}\ a=3x \text{ or } a=2x \qquad  \textbf{(E)}\ a=2x$


 
 \textbf{Answer: }B



\textbf{Solution}

We can factor $a^3-b^3=19x^3$ into:
 \[(a-b)(a^2+ab+b^2)=19x^3.\]Substituting yields:
 \[x(a^2+ab+b^2)=19x^3\]
 \[a^2+ab+b^2=19x^2.\]This is equal to:
 \[(a-b)^2+3ab=19x^2\]
 \[x^2+3ab=19x^2\]
 \[ab=6x^2.\]Checking with the possible answers, along with $a-b=x$ yields the only answer to be $\boxed{(\textbf{B}): a=3x\text{ or }a=-2x}.$

 


	\item (\textit{AHSME 1975}) The equation $x^6 - 3x^5 - 6x^3 - x + 8$ has 


 $\textbf{(A)}\ \text{no real roots} \qquad  \textbf{(B)}\ \text{exactly two distinct negative roots} \qquad \\ \textbf{(C)}\ \text{exactly one negative root} \qquad  \textbf{(D)}\ \text{no negative roots, but at least one positive root} \qquad \\ \textbf{(E)}\ \text{none of these}$


 
 \textbf{Answer: }D



\textbf{Solution 1 (Really Smart)}

Let $P(x) = x^6 - 3x^5 - 6x^3 - x + 8$. When $x < 0$, $P(x) > 0$. Therefore, there are no negative roots. 

 Notice that $P(1) = -1$ and $P(0) = 8$. There must be at least one positive root between 0 and 1, therefore the answer is $\boxed{\textbf{(D)}}$.

 


 
Sidenote : If anyone is wondering what this method is called,

 \begin{verbatim}
It is called the $\textbf {squeeze theorem}$
\end{verbatim}




	\item (\textit{AHSME 1975}) If the $whatsis$ is $so$ when the $whosis$ is $is$ and the $so$ and $so$ is $is \cdot so$, 
what is the $whosis \cdot whatsis$ when the $whosis$ is $so$, the $so$ and $so$ is $so \cdot so$, and the $is$ is two 
($whatsis, whosis, is$, and $so$ are variables taking positive values)? 


 $\textbf{(A)}\ whosis \cdot is \cdot so \qquad  \textbf{(B)}\ whosis \qquad  \textbf{(C)}\ is \qquad  \textbf{(D)}\ so\qquad \textbf{(E)}\ so\text{ and }so$


 
 \textbf{Answer: }E



\textbf{Solution}

From the problem, we are given:

  $whatsis = so$ 

  $whosis = is$ 

  $so + so = is \cdot so$ 

 We want to find what $whatsis \cdot whosis$ is when $whosis = so$, $so + so = so \cdot so$, and $is = 2$.

 Since $is = 2$ and $so + so = is \cdot so = so \cdot so$ (from the given), that means $so = 2$. Now we know that $whatsis = 2$ and $whosis = 2$ as well because they are equal to $so$ and $is$, respectively.

 Plugging in the values, we get $whatsis \cdot whosis = so \cdot is = 4$. This also means that $whatsis \cdot whosis = so + so$.

 Taking a look at the answer choices, we see that our answer is $\boxed{\textbf{(E)}\ so\text{ and }so}$. ~\href{https://artofproblemsolving.com/wiki/index.php/User:Jiang147369}{jiang147369}

 


	\item (\textit{AHSME 1975}) In the sequence of numbers $1, 3, 2, \ldots$ each term after the first two is equal to the term preceding it minus the term preceding that. 
The sum of the first one hundred terms of the sequence is 


 $\textbf{(A)}\ 5 \qquad  \textbf{(B)}\ 4 \qquad  \textbf{(C)}\ 2 \qquad  \textbf{(D)}\ 1 \qquad  \textbf{(E)}\ -1$


 
 \textbf{Answer: }A



\textbf{Solution}

First, write a few terms of the sequence:$1, 3, 2, -1, -3, -2, 1, 3, 2, \ldots$  Notice how the pattern repeats every six terms and every six terms have a sum of 0. Then, find that the $16*6=96$th term is $-2$ and the sum of the all those previous terms is $0$. Then, write the 97th to the 100th terms down: $1, 3, 2,-1$ and add them up to get the sum of $\boxed{\textbf{(A) } 5}$

 


	\item (\textit{AHSME 1975}) If the first term of an infinite geometric series is a positive integer, the common ratio is the reciprocal of a positive integer, 
and the sum of the series is $3$, then the sum of the first two terms of the series is 


 $\textbf{(A)}\ 1/3 \qquad  \textbf{(B)}\ 2/3 \qquad  \textbf{(C)}\ 8/3 \qquad  \textbf{(D)}\ 2 \qquad  \textbf{(E)}\ 9/2$


 
 \textbf{Answer: }C



\textbf{Solution}

Let $a$ and $r$ be the first term and common ratio of the series, respectively.  Then $r = \frac{1}{b}$ for some positive integer $b$.

 The sum of the series is $\frac{a}{1-r} = \frac{a}{1-\frac{1}{b}} = \frac{ab}{b-1} = 3$.

 So $ab = 3(b-1) = 3b - 3$, and thus $ab - 3b = b(a-3) = -3$.

 Since $b$ and $a-3$ are both integers, $b$ is a factor of $3$.  So either $b = 1$ or $b = 3$.  But if $b = 1$, then $a = 0$ which is not positive.  So $b = 3$ and thus $a = 2$.

 Therefore, $a = 2$ and $r = \frac{1}{3}$.  So the sum of the first $2$ terms of the series is $2 + \frac{2}{3} = \boxed{\textbf{(C)}\ \frac{8}{3}}$.

 


	\item (\textit{AHSME 1975}) A man can commute either by train or by bus. If he goes to work on the train in the morning, 
he comes home on the bus in the afternoon; and if he comes home in the afternoon on the train, 
he took the bus in the morning. During a total of $x$ working days, the man took the bus to work 
in the morning $8$ times, came home by bus in the afternoon $15$ times, and commuted by train 
(either morning or afternoon) $9$ times. Find $x$. 


 $\textbf{(A)}\ 19 \qquad  \textbf{(B)}\ 18 \qquad  \textbf{(C)}\ 17 \qquad  \textbf{(D)}\ 16 \qquad \\ \textbf{(E)}\ \text{not enough information given to solve the problem}$


 
 \textbf{Answer: }D



\textbf{Solution 1}

Each day the man takes $2$ trips: a trip to work and a trip back home.  Since the man took $8 + 15 + 9 = 32$ trips, he worked $32 \div 2 = \boxed{\textbf{(D)}\ 16}$ days.

 



\textbf{Solution 2}

The man has three possible combinations of transportation:
 \[\text{Morning train, Afternoon bus (m.t., a.b.)}\]
 \[\text{Morning bus, Afternoon train (m.b., a.t.)}\]
 \[\text{Morning bus, Afternoon bus (m.b, a.b.)}\].

 
Let $y$ be the number of times the man takes the $\text{a.t.}$. Then, $9-y$ is the number of times he takes the $\text{m.t.}$. Keep in mind that $\text{m.b.}=y$ and $\text{a.b.}=9-y$.

 
Let $z$ be the number of times the man takes the $\text{m.b.}$ and $\text{a.b.}$. Now, we get the two equations 
 \[y+z=8\] and 
 \[9-y+z=15.\]

 Solving the system of equations, we get $y=1$ and $z=7$.

 
So during the $x$ working days, the man took the $\text{(m.t., a.b.)}$ on $9-1=8$ days, the $\text{(m.b., a.t.)}$ on $1$ day, and the $\text{(m.b., a.b.)}$ on  $7$ days.

 
Therefore, $x=8+1+7= \boxed{\textbf{(D)}\ 16}$. ~\href{https://artofproblemsolving.com/wiki/index.php?title=User:Jiang147369}{jiang147369}

 


 


	\item (\textit{AHSME 1975}) A positive integer $N$ with three digits in its base ten representation is chosen at random, 
with each three digit number having an equal chance of being chosen. The probability that $\log_2 N$ is an integer is 


 $\textbf{(A)}\ 0 \qquad  \textbf{(B)}\ 3/899 \qquad  \textbf{(C)}\ 1/225 \qquad  \textbf{(D)}\ 1/300 \qquad  \textbf{(E)}\ 1/450$


 
 \textbf{Answer: }D



\textbf{Solution}

In order for $\log_2 N$ to be an integer, $N$ has to be an integer power of $2$. Since $N$ is a three-digit number in decimal form, $N$ can only be $2^{7} = 128,\ 2^{8} = 256,\ \text{or}\ 2^{9} = 512$. From $100$ to $999$, there are $900$ three-digit numbers in total.

 The probability of choosing $128$, $256$, or $512$ from $900$ numbers is $\frac{3}{900} = \frac{1}{300}$, which gives us the answer $\boxed{\textbf{(D)}\ 1/300}$. ~\href{https://artofproblemsolving.com/wiki/index.php/User:Jiang147369}{jiang147369}

 


	\item (\textit{AHSME 1975}) Which positive numbers $x$ satisfy the equation $(\log_3x)(\log_x5)=\log_35$? 


 $\textbf{(A)}\ 3 \text{ and } 5 \text{ only} \qquad  \textbf{(B)}\ 3, 5, \text{ and } 15 \text{ only} \qquad \\ \textbf{(C)}\ \text{only numbers of the form } 5^n \cdot 3^m, \text{ where } n \text{ and } m \text{ are positive integers} \qquad \\ \textbf{(D)}\ \text{all positive } x \neq 1 \qquad  \textbf{(E)}\ \text{none of these}$


 
 \textbf{Answer: }D



\textbf{Solution}

Note that $\log_3 x$ is defined if and only if $x$ is positive, while $\log_x 5$ is defined if and only if $x$ is positive but not equal to $1$.  

 If $x$ is positive and not equal to $1$, by the change of base formula, $(\log_3 x)(\log_x 5) = \frac{\log x}{\log 3}\cdot\frac{\log 5}{\log x} = \frac{\log 5}{\log 3} = \log_3 5$.

 Thus the set of all positive numbers $x$ satisfying the equation is $\boxed{\textbf{(D)}\ \text{all positive } x \neq 1}$.

 


	\item (\textit{AHSME 1975}) In the adjoining figure $\triangle ABC$ is such that $AB = 4$ and $AC = 8$. If $M$ is the midpoint of $BC$ and $AM = 3$, what is the length of $BC$?


 


                    \begin{figure}[H]
                        centering
                        \includegraphics[width=0.5\textwidth]{Problem/AHSME_I_1975_20_Problem_Figure_1.png}
                    \end{figure}
                    


 $\textbf{(A)}\ 2\sqrt{26} \qquad  \textbf{(B)}\ 2\sqrt{31} \qquad  \textbf{(C)}\ 9 \qquad  \textbf{(D)}\ 4+2\sqrt{13} \qquad\\ \textbf{(E)}\ \text{not enough information given to solve the problem}$


 



 
 \textbf{Answer: }B



\textbf{Solution}

Let $BM=CM=x$. Then, by Stewart's Theorem, we have
 \[2x^3+18x=16x+64x\]
 \[\implies x^2+9=40\]
 \[\implies x=\sqrt{31}\implies BC=\boxed{2\sqrt{31}}.\]The answer is $\boxed{B}.$-brainiacmaniac31

 


	\item (\textit{AHSME 1975}) Suppose $f(x)$ is defined for all real numbers $x$; $f(x)>0$ for all $x$, and $f(a)f(b)=f(a+b)$ for all $a$ and $b$. 
Which of the following statements is true? 


 $\text{I. } f(0)=1$


 $\text{II. } f(-a)=1/f(a) \text{ for all } a$


 $\text{III. } f(a) = \sqrt[3]{f(3a)} \text{ for all } a$


 $\text{IV. } f(b)>f(a) \text{ if } b>a$


 

$\textbf{(A)}\ \text{III and IV only} \qquad  \textbf{(B)}\ \text{I, III, and IV only} \qquad  \textbf{(C)}\ \text{I, II, and IV only} \qquad  \textbf{(D)}\ \text{I, II, and III only} \qquad  \textbf{(E)}\ \text{All are true.}$


 
 \textbf{Answer: }D



\textbf{Solution}

$I: f(0) = 1$

 Let $b = 0$. Our equation becomes $f(a)f(0) = f(a)$, so $f(0) = 1$. Therefore $I$ is always true.

 
$II: f(-a) = \frac{1}{f(a)} \text{ for all } a$

 Let $b = -a$. Our equation becomes $f(a)f(-a) = f(0) = 1 \longrightarrow f(-a) = \frac{1}{f(a)}$. Therefore $II$ is always true.

 
$III: f(a) = \sqrt[3] {f(3a)} \text{ for all } a$

 First let $b = a$. We get $f(a)f(a) = f(2a)$. Now let $b = 2a$, giving us $f(3a) = f(a)f(2a) = f(a)^3 \longrightarrow f(a) = \sqrt[3] {f(3a)}$. Therefore $III$ is always true.

 
$IV: f(b) > f(a) \text{ if } b > a$This is false. Let $f(x) = 2^{-x}$, for example. It satisfies the conditions but makes $IV$ false. Therefore $IV$ is not always true.

 
Since $I, II, \text{ and } III$ are true, the answer is $\boxed{\textbf{(D)}}$.

 - mako17

 Note: For $IV$, one may also disprove it with simply $f(x) = 1$.

 


	\item (\textit{AHSME 1975}) If $p$ and $q$ are primes and $x^2 - px + q = 0$ has distinct positive integral roots, 
then which of the following statements are true? 


 $I.\ \text{The difference of the roots is odd.} \\ II.\ \text{At least one root is prime.} \\ III.\ p^2-q\ \text{is prime}. \\ IV.\ p+q\ \text{is prime}$


 

$\textbf{(A)}\ \text{I only} \qquad  \textbf{(B)}\ \text{II only} \qquad  \textbf{(C)}\ \text{II and III only} \qquad  \textbf{(D)}\ \text{I, II and IV only} \qquad  \textbf{(E)}\ \text{All are true.}$


 
 \textbf{Answer: }E



\textbf{Solution}

Since the roots are both positive integers, we can say that $x^2-px+q=(x-1)(x-q)$ since $q$ only has $2$ divisors. Thus, the roots are $1$ and $q$ and $p=q+1$. The only two primes which differ by $1$ are $2,3$ so $p=3$ and $q=2$. $I$ is true because $3-2=1$.$II$ is true because one of the roots is $2$ which is prime.$III$ is true because $3^2-2=7$ is prime.$IV$ is true because $2+3=5$ is prime. Thus, the answer is $\textbf{(E)}$.-brainiacmaniac31

 


	\item (\textit{AHSME 1975}) In the adjoining figure $AB$ and $BC$ are adjacent sides of square $ABCD$; $M$ is the midpoint of $AB$; 
$N$ is the midpoint of $BC$; and $AN$ and $CM$ intersect at $O$. The ratio of the area of $AOCD$ to the area of $ABCD$ is


 
                    \begin{figure}[H]
                        centering
                        \includegraphics[width=0.5\textwidth]{Problem/AHSME_I_1975_23_Problem_Figure_1.png}
                    \end{figure}
                    


 $\textbf{(A)}\ \frac{5}{6}\qquad \textbf{(B)}\ \frac{3}{4}\qquad \textbf{(C)}\ \frac{2}{3}\qquad \textbf{(D)}\ \frac{\sqrt{3}}{2}\qquad \textbf{(E)}\ \frac{(\sqrt{3}-1)}{2}$


 
 \textbf{Answer: }C



\textbf{Solution}

First, let's draw a few auxiliary lines. Drop altitudes from $O$ to $AB$ and from $O$ to $BC$. We can label the points as $J$ and $K$, respectively. This forms square $OKBJ$. Connect $OB$.

 Without loss of generality, set the side length of the square $ABCD$ equal to $1$. Let $NK=x$, and since $N$ is the midpoint of $CB$, $KB$ would be $\frac{1}{2} - x$. With the same reasoning, $MJ=x$ and $JB = \frac{1}{2} - x$

 
                    \begin{figure}[H]
                        centering
                        \includegraphics[width=0.5\textwidth]{Solution/AHSME_I_1975_23_Solution_Figure_1.png}
                    \end{figure}
                    

 We can also see that $\triangle COK$ is similar to $\triangle CMB$. That means $\frac{CK}{OK} = \frac{CB}{MB}$.

 Plugging in the values, we get: $\frac{\frac{1}{2} + x}{\frac{1}{2} - x} = \frac{1}{\frac{1}{2}}$. Solving, we find that $x=\frac{1}{6}$. Then, $OK=\frac{1}{2} - \frac{1}{6} = \frac{1}{3}$. The area of $\triangle COB$ and $\triangle AOB$ together would be $2 (\frac{1}{2} \cdot 1 \cdot \frac{1}{3}) = \frac{1}{3}$. Subtract this area from the total area of $ABCD$ to get the area of $AOCD$.

 So, $AOCD = 1 - \frac{1}{3} = \frac{2}{3}$. The question asks for the ratio of the area of $AOCD$ to the area of $ABCD$, which is $\frac{\frac{2}{3}}{1} = \frac{2}{3}$.

 The answer is $\boxed{\textbf{(C)}\ \frac{2}{3}}$. ~\href{https://artofproblemsolving.com/wiki/index.php/User:Jiang147369}{jiang147369}

 


	\item (\textit{AHSME 1975}) In $\triangle ABC$, $\measuredangle C=\theta$ and $\measuredangle B=2\theta$, where $0^{\circ} <\theta < 60^{\circ}$. 
The circle with center $A$ and radius $AB$ intersects $AC$ at $D$ and intersects $BC$, 
extended if necessary, at $B$ and at $E$ ($E$ may coincide with $B$). Then $EC=AD$


 $\textbf{(A)}\ \text{for no values of}\ \theta \qquad \textbf{(B)}\ \text{only if}\ \theta=45^{\circ} \qquad \textbf{(C)}\ \text{only if}\ 0^{\circ}<\theta\le 45^{\circ}\\ \qquad\ \textbf{(D)}\ \text{only if}\ 45^{\circ}\le\theta < 60^{\circ}\qquad \textbf{(E)}\ \text{for all}\ \theta\ \text{such that}\ 0^{\circ}<\theta < 60^{\circ}$


 
 \textbf{Answer: }E



\textbf{Solution}

Since $AD = AE$, we know $EC = AD$ if and only if triangle $ACE$ is isosceles and $\angle ACE = \angle CAE$. Letting $\angle ACE = \theta$, we want to find when $\angle CAE = \theta$. We know $\angle ABC = \angle AEB = 2\theta$, so $\angle EAB = 180-4\theta$. We also know $\angle CAB = 180-3\theta$, and since $\angle CAE = \angle CAB - \angle EAB$, $\angle CAE = \theta$. Since we now know that $\angle CAE = \angle ACE$ regardless of $\theta$, we have $0^{\circ} < \theta < 60^{\circ}$, or $\boxed{E}$.

 


	\item (\textit{AHSME 1975}) A woman, her brother, her son and her daughter are chess players (all relations by birth). The worst player's twin (who is one of the four players) and the best player are of opposite sex. The worst player and the best player are the same age. Who is the worst player?


 $\textbf{(A)}\ \text{the woman} \qquad \textbf{(B)}\ \text{her son} \qquad \textbf{(C)}\ \text{her brother} \qquad \textbf{(D)}\ \text{her daughter}\\ \qquad \textbf{(E)}\ \text{No solution is consistent with the given information}$


 
 \textbf{Answer: }B



\textbf{Solution 1}

Since the worst player, the worst player's twin, and the best player are all different people of the same age, $3$ of the $4$ players are the same age.  Since the woman is older than her $2$ children, she is not one of these $3$ players.  Therefore, her brother and her $2$ children are the same age.

 So the worst player and worst player's twin must be the woman's $2$ children, while the best player must be the woman's brother.  Therefore, the worst player's twin is the woman's daughter while the worst player is $\boxed{\textbf{(B)}\ \text{her son}}$.

 -j314andrews

 



\textbf{Solution 2}

We know that the worst player's twin and the best player are of opposite sex, and the worst and best players are the same age.

 \textbf{Case 1:} Suppose the daughter was the worst player. The son would have to be her twin because nobody else can be her twin. Then, that means the son and the best player are of opposite sex. But we initially made the daughter the worst player, so she cannot also be the best player. That only leaves the woman to be the best player. However, the daughter and the woman are not the same age. Therefore, the worst player cannot be the daughter.

 \textbf{Case 2:} Suppose the woman was the worst player. Suppose again that her and her brother are twins. Her brother and the best player are of opposite sex. Using the same logic as in Case 1, the best player would be the daughter. The woman and the daughter are not the same age, so the worst player cannot be the woman.

 \textbf{Case 3:} Suppose the brother was the worst player and that him and the woman are twins. The woman and the best player are of opposite sex, so the best player would be the son. If the brother is the woman's twin (our initial assumption), then the son is not the same age as the brother, which means that the son is not the best player, and the brother is not the worst player.

 However, the woman and the brother do not necessarily have to be twins. For this case, we had to assume they were twins or else there was no way to continue the problem, since the brother would have no twin.

 \textbf{Case 4:} Suppose the son was the worst player. The daughter and the best player are of opposite sex, so the best player is the brother. The woman and her brother do not have to be twins. So it is, in fact, possible that the brother and the son are the same age.

 From this, we conclude our answer as $\boxed{\textbf{(B)}\ \text{her son}}$. ~\href{https://artofproblemsolving.com/wiki/index.php/User:Jiang147369}{jiang147369}

 


	\item (\textit{AHSME 1975}) In acute $\triangle ABC$ the bisector of $\measuredangle A$ meets side $BC$ at $D$. The circle with center $B$ and radius $BD$ 
intersects side $AB$ at $M$; and the circle with center $C$ and radius $CD$ intersects side $AC$ at $N$. Then it is always true that


 $\textbf{(A)}\ \measuredangle CND+\measuredangle BMD-\measuredangle DAC =120^{\circ} \qquad \textbf{(B)}\ AMDN\ \text{is a trapezoid}\qquad \textbf{(C)}\ BC\ \text{is parallel to}\ MN\\ \qquad \textbf{(D)}\ AM-AN=\frac{3(DB-DC)}{2}\qquad \textbf{(E)}\ AB-AC=\frac{3(DB-DC)}{2}$


 
 \textbf{Answer: }C



\textbf{Solution(Very cheesy and sketch)}

First by the angle bisector theorem $BM=BD=\frac{AB\cdot BC}{AC+BC}=AB\cdot \frac{BC}{AB+AC}.$ Similarly, $CN=AC\cdot \frac{BC}{AB+AC}.$ We thus have that $AM=AB-AB\cdot \frac{BC}{AB+AC}=AB\cdot \frac{AB+AC-BC}{AB+AC}$ and $AN=AC-CN=AC\cdot \frac{AB+AC-BC}{AB+AC}.$ Finally, since $A,M,B$ and $A,N,C$ are all collinear, there exists a homothety between triangles $BAC$ and $MAN$ with scale factor $\frac{AB+AC-BC}{AB+AC}.$ Thus, $BC\parallel MN,$ or $\boxed{(C)}.$

 --AndrewZhong2012--

 


	\item (\textit{AHSME 1975}) If $p, q$ and $r$ are distinct roots of $x^3-x^2+x-2=0$, then $p^3+q^3+r^3$ equals


 $\textbf{(A)}\ -1 \qquad \textbf{(B)}\ 1 \qquad \textbf{(C)}\ 3 \qquad \textbf{(D)}\ 5 \qquad \textbf{(E)}\ \text{none of these}$


 
 \textbf{Answer: }E



\textbf{Solution 1}

If $p$ is a root of $x^3 - x^2 + x - 2 = 0$, then $p^3 - p^2 + p - 2 = 0$, or
 \[p^3 = p^2 - p + 2.\]Similarly, $q^3 = q^2 - q + 2$, and $r^3 = r^2 - r + 2$, so
 \[p^3 + q^3 + r^3 = (p^2 + q^2 + r^2) - (p + q + r) + 6.\]

 By \href{https://artofproblemsolving.com/wiki/index.php?title=Vieta\%27s\_formulas}{Vieta's formulas}, $p + q + r = 1$, $pq + pr + qr = 1$, and $pqr = 2$. Squaring the equation $p + q + r = 1$, we get
 \[p^2 + q^2 + r^2 + 2pq + 2pr + 2qr = 1.\]Subtracting $2pq + 2pr + 2qr = 2$, we get
 \[p^2 + q^2 + r^2 = -1.\]

 Therefore, $p^3 + q^3 + r^3 = (p^2 + q^2 + r^2) - (p + q + r) + 6 = (-1) - 1  + 6 = \boxed{4}$. The answer is $\text{(E)}$.

 



\textbf{Solution 2(Faster)}

We know that $p^3+q^3+r^3=(p+q+r)(p^2+q^2+r^2-pq-qr-pr)+3pqr$. By Vieta's formulas, $p+q+r=1$,$pqr=2$, and $pq+qr+pr=1$.So if we can find $p^2+q^2+r^2$, we are done. Notice that $(p+q+r)^2=p^2+q^2+r^2+2pq+2qr+2pr$, so $p^2+q^2+r^2=(p+q+r)^2-2(pq+qr+pr)=1^2-2\cdot1=-1$, which means that $p^3+q^3+r^3=1\cdot-2+3\cdot2=\boxed{\text{(E) }4}$ 

 



\textbf{Solution 3 (Beginner's Solution)}

Use Vieta's formulas to get $p+q+r=1$, $pq+qr+pr=1$, and $pqr=2$.

 Square $p+q+r=1$, and get $p^2+q^2+r^2+2pq+2pr+2qr=1$

 Substitute $pq+qr+pr=1$ and simplify to get $-1=p^2+q^2+r^2$

 After that, multiply both sides by $1=p+q+r$, to get $-1=p^3+q^3+r^3+p^2q+q^2r+p^2r+q^2r+r^2p+r^2q$

 Then, factor out $pq$, $qr$, and $pr$:$-1=p^3+q^3+r^3+pq(p+q)+qr(q+r)+pr(p+r)$

 Then, substitute the first equation into $p+q$, $q+r$, and $p+r$.$-1=p^3+q^3+r^3+pq(1-r)+qr(1-p)+pr(1-q)$

 Then, multiply it out:$-1=p^3+q^3+r^3+pq+qr+pr-3pqr$

 After that, substitute the equations $pq+qr+pr=1$ and $pqr=2$:$-1=p^3+q^3+r^3+1-6$

 Solving that, you get $p^3+q^3+r^3=\boxed{\text{(E) }4}$~EZ PZ Ms.Lemon SQUEEZY

 


	\item (\textit{AHSME 1975}) In $\triangle ABC$ shown in the adjoining figure, $M$ is the midpoint of side $BC, AB=12$ and $AC=16$. Points $E$ and $F$ are taken on $AC$ 
and $AB$, respectively, and lines $EF$ and $AM$ intersect at $G$. If $AE=2AF$ then $\frac{EG}{GF}$ equals


 
                    \begin{figure}[H]
                        centering
                        \includegraphics[width=0.5\textwidth]{Problem/AHSME_I_1975_28_Problem_Figure_1.png}
                    \end{figure}
                    


 $\textbf{(A)}\ \frac{3}{2} \qquad \textbf{(B)}\ \frac{4}{3} \qquad \textbf{(C)}\ \frac{5}{4} \qquad \textbf{(D)}\ \frac{6}{5}\\ \qquad \textbf{(E)}\ \text{not enough information to solve the problem}$


 
 \textbf{Answer: }A



\textbf{Solution}

Here, we use Mass Points.Let $AF = x$. We then have $AE = 2x$, $EC = 16-2x$, and $FB = 12 - x$Let $B$ have a mass of $2$. Since $M$ is the midpoint, $C$ also has a mass of $2$. Looking at segment $AB$, we have
 \[2 \cdot (12-x) = \text{m}A_{AB} \cdot x\]So
 \[\text{m}A_{AB} = \frac{24-2x}{x}\]Looking at segment $AC$,we have
 \[2 \cdot (16-2x) = \text{m}A_{AC} \cdot 2x\]So
 \[\text{m}A_{AC} = \frac{16-2x}{x}\]From this, we get
 \[\text{m}E = \frac{16-2x}{x} + 2 \Rrightarrow \text{m}E = \frac{16}{x}\]and
 \[\text{m}F = \frac{24-2x}{x} + 2 \Rrightarrow \text{m}F = \frac{24}{x}\]We want the value of $\frac{EG}{GF}$. This can be written as
 \[\frac{EG}{GF} = \frac{\text{m}F}{\text{m}E}\]Thus
 \[\frac{\text{m}F}{\text{m}E} = \frac {\frac{24}{x}}{\frac{16}{x}} = \frac{3}{2}\]$\boxed{A}$

 



\textbf{Solution 2}

Since we only care about a ratio $EG/GF$, and since we are given $M$ being the midpoint of $\overline{BC}$, we realize we can conveniently also choose $E$ to be the midpoint of $\overline{AC}$. (we're free to choose any point $E$ on $\overline{AC}$ as long as $AE$ is twice $AF$, the constraint given in the problem). This means $AE = 16/2 = 8$, and $AF = 4$. We then connect $ME$ which creates similar triangles $\triangle EMC$ and $\triangle ABC$ by SAS, and thus generates parallel lines $\overline{EM}$ and $\overline{AB}$. This also immediately gives us similar triangles $\triangle AFG \sim \triangle MEG, => EG/FG = ME/AF = 6/4 = \boxed{3/2}$ (note that $ME = 6$ because $ME/AB$ is in $1:2$ ratio).

 


 



\textbf{Solution 3}

In order to find $\frac{EG}{FG}$, we can apply the law of sine to this model. Let:

 $\angle MAC=\alpha, \; \angle MAB=\beta, \; \angle AGE=\gamma,\; \angle AMC=\delta$

 Then, in the $\triangle AMC$ and $\triangle AMB$:

 $\frac{MC}{sin\alpha}=\frac{AC}{sin\delta};$$\frac{MB}{sin\beta}=\frac{AB}{sin\delta}=\frac{MC}{sin\alpha}\cdot\frac{3}{4};$$\frac{sin\alpha}{sin\beta}=\frac{3}{4}$

 In the $\triangle AGE$ and $\triangle AGF$:

 $\frac{FG}{sin\beta}=\frac{AF}{sin\gamma};$$\frac{EG}{sin\alpha}=\frac{AE}{sin\gamma};$$\frac{2FG}{sin\beta}=\frac{EG}{sin\alpha};$$\frac{EG}{FG}=\frac{2sin\alpha}{sin\beta}=\frac{3}{2}$

 Hence, our answer is A.

 -VSN

 


	\item (\textit{AHSME 1975}) What is the smallest integer larger than $(\sqrt{3}+\sqrt{2})^6$?


 $\textbf{(A)}\ 972 \qquad \textbf{(B)}\ 971 \qquad \textbf{(C)}\ 970 \qquad \textbf{(D)}\ 969 \qquad \textbf{(E)}\ 968$


 
 \textbf{Answer: }C



\textbf{Solution(Brute Force)}

$(\sqrt{3}+\sqrt{2})^6=(5+2\sqrt{6})^3=(5+2\sqrt{6})(49+20\sqrt{6})=(485+198\sqrt{6})$ Then, find that $\sqrt{6}$ is about $2.449$. Finally, multiply and add to find that the smallest integer higher is $\boxed {\textbf{(C) } 970}$

 



\textbf{Solution(The Real Way - Binomial Thereom)}

Let's evaluate $(\sqrt{3}+\sqrt{2})^6 + (\sqrt{3}-\sqrt{2})^6$. We see that all the irrational terms cancel. Then, using binomial theorem, we evaluate all the rational terms in the first expression to get 485. Then, the sum of the rational parts of the 2nd term will be 485 as well. Then, we get a total of 970 and since $(\sqrt{3}-\sqrt{2})^6) < 1$, the greatest integer greater than our original expression is $\boxed {\textbf{(C) } 970}$

 


	\item (\textit{AHSME 1975}) Let $x=\cos 36^{\circ} - \cos 72^{\circ}$. Then $x$ equals


 $\textbf{(A)}\ \frac{1}{3}\qquad \textbf{(B)}\ \frac{1}{2} \qquad \textbf{(C)}\ 3-\sqrt{6} \qquad \textbf{(D)}\ 2\sqrt{3}-3\qquad \textbf{(E)}\ \text{none of these}$


 
 \textbf{Answer: }B



\textbf{Solution 1}

Using the difference to product identity, we find that$x=\cos 36^{\circ} - \cos 72^{\circ}$ is equivalent to 
 \[x=\text{-}2\sin{\frac{(36^{\circ}+72^{\circ})}{2}}\sin{\frac{(36^{\circ}-72^{\circ})}{2}} \implies\]
 \[x=\text{-}2\sin54^{\circ}\sin(\text{-}18^{\circ}).\] Since sine is an odd function, we find that $\sin{(\text{-}18^{\circ})}= \text{-} \sin{18^{\circ}}$, and thus $\text{-}2\sin54^{\circ}\sin(\text{-}18^{\circ})=2\sin54^{\circ}\sin18^{\circ}$. Using the property $\sin{(90^{\circ}-a)}=\cos{a}$, we find
 \[x=2\cos(90^{\circ}-54^{\circ})\cos(90^{\circ}-18^{\circ}) \implies\]
 \[x=2\cos36^{\circ}\cos72^{\circ}.\]We multiply the entire expression by $\sin36^{\circ}$ and use the double angle identity of sine twice to find
 \[x\sin36^{\circ}=2\sin36^{\circ}\cos36^{\circ}\cos72^{\circ} \implies\]
 \[x\sin36^{\circ}=\sin72^{\circ}\cos72^{\circ} \implies\]
 \[x\sin36^{\circ}=\frac{1}{2}\sin144^{\circ}.\]Using the property $\sin(180^{\circ}-a)=\sin{a}$, we find $\sin144^{\circ}=\sin36^{\circ}.$ Substituting this back into the equation, we have
 \[x\sin36^{\circ}=\frac{1}{2}\sin36^{\circ}.\]Dividing both sides by $\sin36^{\circ}$, we have
 \[x=\boxed{\textbf{(B)}\ \frac{1}{2}}\]

 



\textbf{Solution 2}

We observe that $72^{\circ}$ is twice of $36^{\circ}$, so evaluating $\cos 36^{\circ}$ and using the double angle identify is enough to find $\cos 36^{\circ} - \cos 72^{\circ}$ 

 We first let $n = \cos 36^{\circ}$. Notice that $36^{\circ}/4 = 9^{\circ}$ and $36^{\circ} = 45^{\circ} - 9^{\circ}$. Rewriting $\cos 36^{\circ}$ as $\cos (45^{\circ} - 9^{\circ})$, we then use the difference of angles identify to find 
 \[n = \cos 36^{\circ} = \cos (45^{\circ} - 9^{\circ}) = \cos 45^{\circ}\cos 9^{\circ} - \sin 45^{\circ}\sin 9^{\circ} = \frac{\sqrt2}{2}(\cos 9^{\circ} - \sin 9^{\circ})\] Rearranged to get 
 \[\sqrt2 n = \cos 9^{\circ} - \sin 9^{\circ}\]Using the half-angle identify twice, and noticing that $\cos 9^{\circ}$ is positive, we get 
 \[\cos 9^{\circ} = \cos \frac{18^{\circ}}{2} = \sqrt \frac{1 + \cos 18^{\circ}}{2} = \sqrt \frac{1 + \sqrt \frac{1 + \cos 36^{\circ}}{2}}{2}\]Also, 
 \[\sin 9^{\circ} = \sqrt {1 - \cos^2 9} = \sqrt {1 - \frac{1 + \sqrt{\frac{1 + 36^{\circ}}{2}}}{2}} = \sqrt {\frac{1 - \sqrt{\frac{1 + 36^{\circ}}{2}}}{2}}\]

 Thus we can get
 \[\sqrt2 n = \sqrt \frac{1 + \sqrt \frac{1 + \cos 36^{\circ}}{2}}{2} - \sqrt {\frac{1 - \sqrt{\frac{1 + 36^{\circ}}{2}}}{2}}\]Squaring both sides and simplifying, this results in 
 \[2n^2 = 1 - \sqrt{\frac{1-n}{2}}\]From which we get 
 \[2n^2 - 1 = - \sqrt{\frac{1-n}{2}}\]Squaring both sides again, 
 \[4n^4 - 4n^2 +1 = \frac{1-n}{2}\]Multiplying both sides by $2$ and rearranging, 
 \[8n^4 - 8n^2 + n + 1 = 0\]Factoring out $8n^2$ from the first two terms and applying difference of squares, 
 \[8n^2(n+1)(n-1) + n+1 = 0\] We know that $n = \cos 36^{\circ}$ which is not equal to $-1$, so we divide both sides by $n+1$ 
 \[8n^3 -8n^2 +1 = 0\]Using the rational root test, we find that $n = \frac{1}{2}$ is a solution of this equation. However, we also know that $\cos 36^{\circ}$ is not equal to $\frac{1}{2}$, so we divide both sides by $n- \frac{1}{2}$ to get 
 \[8n^2 - 4n - 2 = 0\] Using the quadratic formula, we find that 
 \[n = \frac{1 \pm \sqrt{5}}{4}\] We know $n$ must be positive, which gives us 
 \[n = \frac{1 + \sqrt{5}}{4}\]

 Thus, using the double angle identify, we can find 
 \[x = \cos 36^{\circ} - \cos 72^{\circ} = \cos 36^{\circ} - 2\cos^2 36^{\circ} + 1 = n - 2n^2 + 1 = \boxed{\textbf{(B)} \frac{1}{2}}\]

 -SaraLiu24

 



\textbf{Solution 3 (Roots of Unity)}

Let $w = \cos 72^\circ - i \sin 72^\circ$.  Then $w$ is a fifth root of unity, so $w^5 = 1$.  So $w^5-1 = (w-1)(w^4 + w^3 + w^2 + w + 1) = 0$.  Since $w - 1 \neq 0$, $w^4 + w^3 + w^2 + w + 1 = 0$.

 Both the real and imaginary parts of $w^4 + w^3 + w^2 + w + 1$ must be equal to $0$, and the real part of $w^4 + w^3 + w^2 + w + 1$ can also be expressed as $\cos 288^\circ + \cos 216^\circ + \cos 144^\circ + \cos 72^\circ + 1$.  Also, $\cos 288^\circ = \cos 72^\circ$ and $\cos 216^\circ = \cos 144^\circ = -\cos 36^\circ$.

 So $\cos 288^\circ + \cos 216^\circ + \cos 144^\circ + \cos 72^\circ + 1 = \cos 72^\circ - \cos 36^\circ - \cos 36^\circ + \cos 72^\circ + 1 = 2(\cos 72^\circ - \cos 36^\circ) + 1 = 0$.  Therefore, $\cos 72^\circ - \cos 36^\circ = -\frac{1}{2}$ and $\cos 36^\circ - \cos 72^\circ = \boxed{\textbf{(B)}\ \frac{1}{2}}$.

 -j314andrews

 


\end{enumerate}